% ============================================================
% SWA-MPPI v3: Cross-Cultural Moral Alignment of LLMs
% NeurIPS 2026 Submission
% Compile: pdflatex -> bibtex -> pdflatex -> pdflatex
% ============================================================
\documentclass{article}
\usepackage[final,main]{neurips_2026}

\usepackage[utf8]{inputenc}
\usepackage[T1]{fontenc}
\usepackage{hyperref}
\usepackage{url}
\usepackage{booktabs}
\usepackage{amsfonts}
\usepackage{amsmath}
\usepackage{amssymb}
\usepackage{nicefrac}
\usepackage{microtype}
\usepackage{xcolor}
\usepackage{graphicx}
\usepackage{subcaption}
\usepackage{algorithm}
\usepackage{algpseudocode}
\usepackage{multirow}
\usepackage{makecell}
\usepackage{bbm}

% Colour macros
\definecolor{improvegreen}{RGB}{0,128,0}
\definecolor{worsered}{RGB}{180,0,0}
\newcommand{\gain}[1]{\textcolor{improvegreen}{#1}}
\newcommand{\loss}[1]{\textcolor{worsered}{#1}}

% ============================================================
\title{%
  \textbf{SWA-MPPI}: Socially-Weighted Alignment with\\
  Model Predictive Path Integral Control\\
  for Cross-Cultural Moral Alignment of Large Language Models%
}

\author{%
  Anonymous Authors\\
  \texttt{anonymous@institution.edu}
}

\begin{document}

\maketitle

% ============================================================
\begin{abstract}
% ============================================================
Large language models (LLMs) encode a culturally monolithic moral stance that
systematically diverges from the diverse ethical preferences of populations worldwide.
We introduce \textbf{SWA-MPPI}~(\textbf{S}ocially-\textbf{W}eighted \textbf{A}lignment
with \textbf{M}odel \textbf{P}redictive \textbf{P}ath \textbf{I}ntegral control), a
training-free, inference-time framework that steers LLM moral decisions toward
country-specific human preferences without modifying model weights. For each target
country, SWA-MPPI constructs a small ensemble of culturally grounded \emph{persona agents}
from World Values Survey (WVS) Wave~7 data, evaluates their logit-space consensus in a
single batched forward pass, and applies a Prospect-Theory-weighted MPPI optimisation
step whenever inter-agent disagreement exceeds an adaptively calibrated
threshold~$\tau$. Evaluated on the Multilingual Trolley Problem (MultiTP) benchmark
across 15 countries, 12 native languages, and 6 moral dimensions, SWA-MPPI reduces
Jensen-Shannon Distance to human moral preferences by \textbf{37.6\%} (Qwen2.5-72B)
and \textbf{34.7\%} (Llama-3.1-70B) on average relative to vanilla baselines, while
improving Pearson correlation with human preference vectors by up to $+0.41$ points.
The framework requires no labelled preference data, operates on frozen weights, and
introduces a mean per-scenario overhead of approximately 730\,ms on a single
NVIDIA H100 GPU.
\end{abstract}

% ============================================================
\section{Introduction}
\label{sec:intro}
% ============================================================

As large language models are deployed across increasingly diverse cultural contexts, a
fundamental tension has emerged between the uniformity of a single trained model and the
plurality of human moral values worldwide. Models trained predominantly on English-language
corpora tend to inherit the ethical intuitions of Western, educated, industrialised, rich,
and democratic (WEIRD) populations~\citep{henrich2010weirdest}, yet they are increasingly
enlisted to assist with consequential decisions in societies whose moral norms differ
substantially~\citep{awad2018moral, jin2025multitp}.

The Moral Machine experiment~\citep{awad2018moral} collated over 40~million human
judgments from 233~countries on trolley-problem scenarios, demonstrating profound
cross-cultural variation in moral priorities: East Asian respondents weight youth more
heavily; Southern-hemisphere populations show distinct utilitarian calculi; Islamic
societies prioritise different social roles. The multilingual successor benchmark,
MultiTP~\citep{jin2025multitp}, operationalises this evidence into country-specific
Average Marginal Component Effects (AMCE) across six moral dimensions and 107 languages,
providing a large-scale, cross-culturally validated ground truth for evaluating LLM moral
alignment. Against this benchmark, \citet{jin2025multitp} found that none of 19 evaluated
contemporary LLMs achieved broadly satisfactory cross-cultural alignment: most exhibit
WEIRD-biased moral reasoning that departs markedly from the graded distributions observed
in non-WEIRD human respondents.

Existing approaches to value alignment fall into two broad families. \emph{Fine-tuning}
on human preference data~\citep{bai2022constitutional, ouyang2022instructgpt} can improve
cultural fit but demands large per-culture labelled datasets, modifies weights irreversibly,
and scales poorly across the 130+ countries in MultiTP. \emph{Prompting-based steering}
is lightweight but typically relies on ad-hoc system instructions with no empirical grounding
in population-level values, and its effects on cross-lingual moral reasoning are
inconsistent~\citep{jin2025multitp}. Activation-level interventions~\citep{arditi2025refusal,
zou2023representation} manipulate internal representations but require white-box access and
careful per-model direction engineering.

We propose \textbf{SWA-MPPI} (Socially-Weighted Alignment with Model Predictive Path
Integral control), a principled, training-free inference-time framework grounded in three
observations. First, moral preferences within any society are not monolithic: age cohorts,
ideological subgroups, and social strata exhibit systematically different moral
priorities~\citep{awad2018moral, kleiman2017learning}. Second, the World Values
Survey~\citep{haerpfer2022wvs} provides rich, empirically calibrated data on these
within-country sub-group differences across 80+ nations. Third, when culturally grounded
persona agents derived from WVS data \emph{disagree} on a decision, that disagreement is
itself an informative signal that can be exploited through stochastic optimal
control~\citep{williams2017mppi} to produce a more culturally representative output.

SWA-MPPI operates as follows. For a target country, it constructs $N=4$ system-prompt
personas from WVS Wave~7 statistics, covering three age cohorts and a utilitarian agent.
It evaluates all $N+1$ agents (including the base model) in a single batched forward pass
and extracts \texttt{LEFT}/\texttt{RIGHT} decision logits. If the inter-agent variance of
decision gaps exceeds an adaptively calibrated threshold~$\tau$, it samples $K=128$
perturbations around the consensus signal and scores each via a Prospect Theory value
function~\citep{kahneman1979prospect} that balances private cultural utility with social
welfare (cooperation weight $\lambda_{\text{coop}}=0.7$). The importance-weighted MPPI
update then yields a minimal-KL correction $\delta^*$ to the consensus gap; otherwise the
consensus is used directly. A positional debiasing step averages two passes with
swapped group orderings to cancel recency effects.

\paragraph{Contributions.}
\begin{itemize}
  \item We introduce SWA-MPPI, a training-free inference-time alignment framework that
        combines WVS-grounded cultural personas, adaptive conflict detection, and
        Prospect-Theory-weighted MPPI optimisation in logit space (\S\ref{sec:method}).
  \item We demonstrate mean JSD reductions of \textbf{37.6\%} (Qwen2.5-72B) and
        \textbf{34.7\%} (Llama-3.1-70B) versus vanilla baselines across 15 countries,
        with gains reaching $+75\%$ for Anglosphere country--language pairs
        (\S\ref{sec:results}).
  \item We show that SWA-MPPI outperforms all prompting baselines and activation
        steering on Llama-3.1-70B, with a mean JSD gap of $+34.7\%$ versus the best
        alternative baseline (\S\ref{sec:results}).
  \item We conduct a systematic ablation isolating the contributions of personas, MPPI,
        and the Prospect Theory utility function, finding that persona removal is the
        single most damaging ablation ($-38.1\%$ Pearson) (\S\ref{sec:ablation}).
\end{itemize}

% ============================================================
\section{Related Work}
\label{sec:related}
% ============================================================

\paragraph{LLM moral alignment benchmarks.}
Several benchmarks probe whether LLMs reason morally in ways that match human judgments.
\citet{hendrycks2021aligning} propose the ETHICS dataset covering commonsense morality,
deontology, and utilitarianism, but in English only with a small crowdsourced ground truth.
The Commonsense Norm Bank~\citep{commonsense2021} provides 1.7M judgments but lacks
systematic parametric variation. MoralExceptQA~\citep{jin2022exceptions} and
MoCa~\citep{nie2023moca} introduce more controlled variation but remain English-centric.
PRISM~\citep{kirk2024prism} and GlobalOpinionQA~\citep{durmus2023global} address
cross-cultural diversity but cover only a handful of languages. The Moral Machine
experiment~\citep{awad2018moral} provides the largest cross-cultural dataset (40M+
judgments, 233~countries); \citet{takemoto2024moral} analyse LLM responses to Moral
Machine scenarios but without cross-linguistic grounding. MultiTP~\citep{jin2025multitp}
is the most comprehensive benchmark, supporting 107~languages and grounding evaluation
in the Moral Machine data---we adopt it as our primary benchmark and extend it by
proposing an alignment \emph{intervention} rather than a diagnostic study alone.

\paragraph{Cross-cultural LLM evaluation.}
A growing body of work documents how LLMs favour particular cultural perspectives.
\citet{santurkar2023opinions} show that LLM outputs correlate most strongly with liberal,
educated American respondents. \citet{durmus2023global} observe degraded alignment with
subjective global opinions for non-English queries. \citet{ryan2024unintended} identify
unintended cultural impacts of RLHF alignment, while \citet{manvi2024geographic}
demonstrate geographic bias in geospatial reasoning. \citet{deshpande2023toxicity} show
that assigning cultural personas to LLMs can surface toxic outputs, motivating careful
grounding of persona construction. \citet{arora2023probing} probe LLMs with cross-cultural
survey questions and find systematic WEIRD bias. The recurring finding across these
studies---that models underserve non-WEIRD populations---motivates the cultural
calibration objective of SWA-MPPI.

\paragraph{Inference-time alignment and steering.}
Methods for steering LLM outputs at inference time without retraining have attracted
growing interest. Activation steering~\citep{arditi2025refusal} and representation
engineering~\citep{zou2023representation} manipulate internal activations but require
white-box model access and direction-specific engineering. In the prompting space,
Constitutional AI~\citep{bai2022constitutional} uses iterative self-critique loops, and
chain-of-thought prompting~\citep{wei2022chain} elicits structured reasoning. Persona
and role-play prompting~\citep{deshpande2023toxicity} uses system-level character
assignments to modulate outputs; however, without empirical grounding such personas
may introduce stereotyping. Closest to our prompt-level baselines, PRISM-style
prompting~\citep{kirk2024prism} and WVS profile prompting provide cultural context in
the system message but without the multi-agent ensemble or control-theoretic optimisation
of SWA-MPPI. \citet{sorensen2024roadmap} articulate the goal of \emph{pluralistic
alignment}---systems that model the diversity of human values rather than collapsing to
a single stance---which SWA-MPPI operationalises concretely.

\paragraph{Model Predictive Path Integral (MPPI) control.}
MPPI~\citep{williams2017mppi} is a sampling-based stochastic optimal control algorithm
that replaces gradient computation with Monte Carlo rollouts of perturbations, weighting
each sample by its exponentiated utility. Originally developed for autonomous driving
and robotics, MPPI has recently been adapted for token-level generation steering in
LLMs~\citep{huang2024mppi}. Our work is the first, to our knowledge, to apply MPPI in
\emph{logit space} for cross-cultural moral alignment, combining it with the
Prospect Theory value function~\citep{kahneman1979prospect} to model asymmetric
gain-loss sensitivity characteristic of human moral judgment.

\paragraph{Prospect Theory in computational models.}
Kahneman and Tversky's Prospect Theory~\citep{kahneman1979prospect} provides a
descriptively accurate account of human decision-making under risk, capturing diminishing
sensitivity ($\alpha=\beta=0.88$) and loss aversion ($\kappa=2.25$). It has been applied
widely in computational moral psychology~\citep{kleiman2017learning} and risk-sensitive
reinforcement learning. We adopt it as the utility backbone of SWA-MPPI's collective
welfare scoring, arguing that human moral preferences exhibit exactly the same asymmetric
curvature that Prospect Theory was designed to capture.

% ============================================================
\section{Method: SWA-MPPI}
\label{sec:method}
% ============================================================

\subsection{Problem Formulation}
\label{sec:method:problem}

Let $\mathcal{X}$ be the set of moral dilemma scenarios in the MultiTP benchmark,
each of which presents an LLM with a forced binary choice encoded as decision
tokens \texttt{LEFT}~($\ell$) and \texttt{RIGHT}~($r$). For a target country $c$,
let $\mathbf{h}^{(c)} \in \mathbb{R}^6$ be the human AMCE preference vector, where
each coordinate $h_d^{(c)}$ encodes the Average Marginal Component Effect of moral
dimension $d \in \{\text{Species, Gender, Age, Fitness, SocialValue, Utilitarianism}\}$
derived from 40M+ human judgments~\citep{awad2018moral}.

A vanilla LLM $f_\theta$ with frozen parameters $\theta$ produces logits
$\mathbf{z}(\mathbf{x}) = [z_\ell,\, z_r] \in \mathbb{R}^2$ at the final token position.
The probability of sparing the preferred group is $p_{\text{spare}} = \sigma(\Delta / T_{\text{dec}})$,
where $\Delta = z_r - z_\ell$ and $T_{\text{dec}}=0.5$ is the decision temperature.

\noindent\textbf{AMCE estimation.}
We follow the evaluation protocol of \citet{jin2025multitp} and estimate each
dimension's AMCE from model outputs as the \emph{marginal preference rate} (MPR)
for that dimension's preferred choice:
\begin{equation}
  \hat{m}_d^{(c)} = \frac{1}{|\mathcal{S}_d^{(c)}|}\sum_{\mathbf{x} \in \mathcal{S}_d^{(c)}}
  p_{\text{spare}}(\mathbf{x}),
  \label{eq:mpr}
\end{equation}
where $\mathcal{S}_d^{(c)}$ is the set of scenarios for dimension~$d$ in country~$c$
and $p_{\text{spare}}(\mathbf{x})$ is the model's predicted probability of sparing the
preferred group in scenario~$\mathbf{x}$.
This estimator is consistent with the MultiTP benchmark's evaluation convention.
The human ground-truth AMCEs in \texttt{country\_specific\_ACME.csv}~\citep{jin2025multitp}
are derived from the same marginal averaging over the 40M+ Moral Machine responses,
ensuring that model and human estimates are computed on a comparable basis.
We note that a full conjoint regression estimator would produce AMCEs corrected for
co-occurrence patterns across dimensions; however, \citet{awad2018moral} show that
marginal averages are unbiased estimators under the balanced randomisation design of
the Moral Machine, making Eq.~\ref{eq:mpr} a valid and reproducible alignment metric.

\textbf{Goal:} Without modifying $\theta$, find an inference-time logit correction
$\delta^*(\mathbf{x}, c)$ such that the corrected AMCE vector
$\hat{\mathbf{m}}^{(c)}(\theta, \delta^*)$ is closer to $\mathbf{h}^{(c)}$
under Jensen-Shannon Divergence (JSD).

\subsection{Cultural Persona Construction from WVS Data}
\label{sec:method:personas}

For each target country $c$, SWA-MPPI constructs $N=4$ system-prompt persona agents
grounded in World Values Survey Wave~7 data~\citep{haerpfer2022wvs}. WVS Wave~7 provides
nationally representative survey data across 8 value dimensions used in persona
construction: \emph{gender egalitarianism}, \emph{religious importance},
\emph{interpersonal trust}, \emph{moral permissiveness}, \emph{work centrality},
\emph{family importance}, \emph{personal autonomy}, and \emph{meritocratic orientation}.

Three personas correspond to age cohorts: \emph{young} (18--35), \emph{middle-aged}
(36--55), and \emph{older} ($>$55). Each cohort persona is parameterised by WVS
statistics (means and direction of within-cohort variation) for the target country.
A fourth, country-invariant persona is a \emph{utilitarian agent} that maximises
the number of lives saved irrespective of the identities of the individuals. This
design encodes both culturally specific variation (the three cohort personas) and a
universal moral reference point (the utilitarian).

Each persona is instantiated as a natural-language system prompt of the form:

\begin{quote}
\small\itshape
You are a [age-cohort] adult from [country]. The cultural values typical of your
social group are: [religiosity], [gender egalitarianism], [personal autonomy], and
[meritocratic orientation]. When making moral judgments, you tend to [value
description reflecting WVS profile].
\end{quote}

\noindent WVS grounding ensures that persona descriptions reflect empirically measured
within-country variation rather than author-imposed stereotypes or uniform national
identities~\citep{deshpande2023toxicity}. Personas are constructed once per country
and shared across all scenarios.

To illustrate, the young-adult persona for the \textbf{United States} reads:
\begin{quote}
\small\itshape
You are a young adult (18--35) from the United States of America. Based on survey
data, people in your age group in the USA typically value: strong gender equality
(score: 8.2/10), low religious centrality in moral decisions (score: 3.1/10), high
personal autonomy (score: 7.8/10), high meritocratic orientation (score: 7.5/10),
moderate family importance (score: 6.4/10), and moderate interpersonal trust
(score: 5.3/10). When making moral judgments, you tend to emphasise individual
rights and equal treatment, and weigh personal achievement and rule-following.
\end{quote}
Contrast this with the older-adult persona for \textbf{Vietnam}:
\begin{quote}
\small\itshape
You are an older adult (55+) from Vietnam. Based on survey data, people in your age
group in Vietnam typically value: moderate gender equality (score: 5.1/10), high
religious and moral traditionalism (score: 7.2/10), strong family importance
(score: 9.4/10), moderate personal autonomy (score: 4.8/10), and low interpersonal
trust toward strangers (score: 3.2/10). When making moral judgments, you tend to
emphasise collective welfare, family obligations, and hierarchical social roles.
\end{quote}
The utilitarian persona is country-invariant and reads: \textit{``You are a strict
utilitarian. Always choose the option that saves the greatest number of lives,
regardless of the characteristics of the individuals involved.''}

\subsection{Batched Agent Evaluation and Logit Extraction}
\label{sec:method:eval}

Given scenario $\mathbf{x}$, SWA-MPPI constructs $N+1=5$ input sequences by prepending
the base system prompt ($\varnothing$) and each of the $N$ persona prefixes
$\{\mathbf{s}_i\}_{i=1}^N$ to the query. All $N+1$ sequences are padded to the same
length and evaluated in a \emph{single} batched forward pass:
\begin{equation}
  \mathbf{z}_{\text{base}},\; \mathbf{z}_1, \ldots, \mathbf{z}_N
  = f_\theta\!\left(\mathbf{x};\; \varnothing,\; \mathbf{s}_1, \ldots, \mathbf{s}_N\right),
  \label{eq:batch}
\end{equation}
where each $\mathbf{z}_i = [z_{i,\ell},\; z_{i,r}]$ is the raw logit pair for tokens
\texttt{LEFT} and \texttt{RIGHT}. A per-category logit temperature $T_{\text{cat}}$
(Table~\ref{tab:hyperparams}) scales logits before computing decision gaps, modulating
the confidence of categories that are inherently less discriminative (e.g., Species
logits tend to be extreme and require softening with $T_{\text{cat}}=4.0$).

Agent decision gaps and inter-agent rewards are defined as:
\begin{align}
  \delta_i &= \frac{z_{i,r} - z_{i,\ell}}{T_{\text{cat}}}, \quad i \in \{1, \ldots, N\},
  \label{eq:gap} \\
  \bar{\delta} &= \frac{1}{N}\sum_{i=1}^{N} \delta_i \quad \text{(consensus signal)},
  \label{eq:consensus} \\
  r_i &= \delta_i - \delta_{\text{base}} \quad \text{(per-agent reward)},
  \label{eq:reward}
\end{align}
where $\delta_{\text{base}} = (z_{\text{base},r} - z_{\text{base},\ell})/T_{\text{cat}}$
is the base-model decision gap. The reward $r_i$ measures how much persona $i$ shifts
the model's preference relative to the culturally unanchored default.

\subsection{Adaptive Conflict Detection}
\label{sec:method:tau}

When persona agents agree---i.e., their decision gaps cluster tightly---the consensus
$\bar{\delta}$ provides a reliable cultural signal, and no further optimisation is
needed. When agents disagree substantially, the vanilla model's output may not reflect
the target country's moral balance, warranting a controlled correction.

SWA-MPPI operationalises this via a per-country threshold $\tau^{(c)}$ calibrated on a
small set of $n_{\text{cal}}=50$ scenarios such that the MPPI \emph{trigger rate}
(fraction of scenarios for which MPPI is activated) approximates a target rate
$\rho_{\text{target}}=0.35$:
\begin{equation}
  \tau^{(c)} = \operatorname{Percentile}_{(1-\rho_{\text{target}})\times 100}
  \!\Bigl(\bigl\{\operatorname{Var}(\delta_1^{(s)}, \ldots, \delta_N^{(s)})\bigr\}_{s=1}^{n_{\text{cal}}}\Bigr).
  \label{eq:tau}
\end{equation}
The MPPI step is triggered on scenario $\mathbf{x}$ if and only if
$\nu \triangleq \operatorname{Var}(\delta_1, \ldots, \delta_N) \geq \tau^{(c)}$.
This adaptive calibration allows the trigger rate to naturally vary across countries
in proportion to the true degree of inter-persona disagreement, rather than being
fixed globally.

\subsection{Prospect-Theory-Weighted MPPI Optimisation}
\label{sec:method:mppi}

When triggered, SWA-MPPI draws $K=128$ perturbation samples
$\epsilon_k \sim \mathcal{N}(0,\sigma^2)$ with $\sigma=0.3$ and evaluates a collective
welfare score for each perturbed decision gap $\tilde{\delta}_k = \bar{\delta} + \epsilon_k$.

\paragraph{Per-agent cooperative utility.}
The utility of perturbation $\epsilon_k$ for persona agent $i$ combines a private utility
based on agent $i$'s own reward with a social utility based on the average reward of the
remaining agents:
\begin{equation}
  U_i(\epsilon_k) = (1 - \lambda_{\text{coop}})\,v(r_i\,\tilde{\delta}_k)
  \;+\; \lambda_{\text{coop}}\,v\!\Bigl(\bar{r}_{-i}\,\tilde{\delta}_k\Bigr),
  \label{eq:util_agent}
\end{equation}
where $\bar{r}_{-i} = \frac{1}{N-1}\sum_{j \neq i} r_j$ is the mean reward of all agents
except $i$, $\lambda_{\text{coop}}=0.7$ is the cooperation weight, and $v(\cdot)$ is the
Prospect Theory value function of \citet{kahneman1979prospect}:
\begin{equation}
  v(x) = \begin{cases}
    x^{\alpha} & x \geq 0,\\
    -\kappa\,|x|^{\beta} & x < 0,
  \end{cases}
  \label{eq:prospect}
\end{equation}
with diminishing sensitivity exponents $\alpha=\beta=0.88$ and loss-aversion coefficient
$\kappa=2.25$. The Prospect Theory parameterisation captures the empirically documented
asymmetry between moral gains and losses: agents penalise culturally misaligned decisions
more steeply than they reward aligned ones, consistent with the loss-aversion observed in
human moral judgment~\citep{kleiman2017learning}.

\paragraph{Collective utility and KL regularisation.}
The total collective utility aggregates across agents and adds a KL penalty that prevents
excessively large deviations from the prior:
\begin{equation}
  U_{\text{total}}(\epsilon_k)
  = \frac{1}{N}\sum_{i=1}^{N} U_i(\epsilon_k)
  \;-\; \alpha_{\text{KL}}\,\frac{\epsilon_k^2}{2\sigma^2},
  \label{eq:util_total}
\end{equation}
where $\alpha_{\text{KL}}=0.05$. The KL term is the log-ratio of the proposal
$\mathcal{N}(\tilde{\delta}_k, \sigma^2)$ to the prior $\mathcal{N}(\bar{\delta}, \sigma^2)$,
which for scalar perturbations reduces to $\epsilon_k^2/(2\sigma^2)$.

\paragraph{MPPI importance weighting.}
Following the MPPI update rule~\citep{williams2017mppi}, importance weights are computed
via a softmax over sample utilities, and the optimal correction is their weighted sum:
\begin{equation}
  w_k = \frac{\exp\!\bigl(U_{\text{total}}(\epsilon_k)/\beta\bigr)}
             {\sum_{k'}\exp\!\bigl(U_{\text{total}}(\epsilon_{k'})/\beta\bigr)},
  \qquad
  \delta^* = \sum_{k=1}^{K} w_k\,\epsilon_k,
  \label{eq:mppi}
\end{equation}
where $\beta = T_{\text{dec}} = 0.5$ is the MPPI temperature, matching the decision
temperature. The final optimal signal is $\delta_{\text{opt}} = \bar{\delta} + \delta^*$.

\subsection{Positional Debiasing}
\label{sec:method:debias}

LLMs exhibit a known tendency to favour options appearing later in a prompt
(\emph{recency bias})~\citep{liu2024lost}. To cancel this effect, SWA-MPPI evaluates
each scenario twice: once with the original group ordering (Group~A in the LEFT lane,
Group~B in the RIGHT lane) and once with the ordering swapped. The decision gaps from
the two passes are averaged before computing the consensus and applying any MPPI
correction. This positional debiasing is applied to both the base model and all persona
agents within the same batched pass, requiring two batched forward passes per scenario in
total.

\subsection{Final Prediction}
\label{sec:method:final}

The final probability of sparing the preferred group is:
\begin{equation}
  p_{\text{spare}}
  = \sigma\!\Bigl(\frac{\delta_{\text{opt}}}{T_{\text{dec}}}\Bigr),
  \quad T_{\text{dec}} = 0.5.
  \label{eq:final}
\end{equation}

Algorithm~\ref{alg:swa} summarises the complete SWA-MPPI inference procedure.
Table~\ref{tab:hyperparams} lists all hyperparameters.

\begin{algorithm}[t]
\caption{SWA-MPPI Inference for Scenario $\mathbf{x}$ in Country $c$}
\label{alg:swa}
\begin{algorithmic}[1]
\Require Frozen LLM $f_\theta$; persona prompts $\{\mathbf{s}_i\}_{i=1}^{N}$;
         calibrated threshold $\tau^{(c)}$; hyperparameters $K,\sigma,\lambda_{\text{coop}},\alpha_{\text{KL}},T_{\text{dec}}$
\Ensure Alignment-corrected probability $p_{\text{spare}} \in [0,1]$

\State \textbf{// Positional debiasing: run original and swapped orderings}
\For{$\text{order} \in \{\text{orig}, \text{swap}\}$}
  \State \textbf{Batched forward:} evaluate $f_\theta$ with base and all $N$ persona prefixes in one batch
  \State Extract logit pairs $\mathbf{z}_{\text{base}}^{(\text{ord})},\,\mathbf{z}_1^{(\text{ord})},\ldots,\mathbf{z}_N^{(\text{ord})}$ \hfill\Comment{Eq.~\ref{eq:batch}}
  \State Compute $\delta_i^{(\text{ord})} = (z_{i,r}^{(\text{ord})} - z_{i,\ell}^{(\text{ord})})/T_{\text{cat}}$ for $i \in \{0,\ldots,N\}$
\EndFor
\State Average: $\delta_i \leftarrow (\delta_i^{(\text{orig})} + (-\delta_i^{(\text{swap})}))/2$ for all $i$ \hfill\Comment{Debiasing}
\State Compute consensus $\bar{\delta}$ and per-agent rewards $r_i$ \hfill\Comment{Eqs.~\ref{eq:consensus}--\ref{eq:reward}}
\State Compute inter-agent variance $\nu = \operatorname{Var}(\delta_1,\ldots,\delta_N)$
\If{$\nu \geq \tau^{(c)}$} \hfill\Comment{Conflict detected; run MPPI}
  \For{$k = 1$ \textbf{to} $K$}
    \State Sample $\epsilon_k \sim \mathcal{N}(0, \sigma^2)$; set $\tilde{\delta}_k = \bar{\delta} + \epsilon_k$
    \State Compute $U_i(\epsilon_k)$ for all $i$ via Eq.~\ref{eq:util_agent} (Prospect Theory)
    \State Compute $U_{\text{total}}(\epsilon_k)$ via Eq.~\ref{eq:util_total} (with KL penalty)
  \EndFor
  \State Compute importance weights $\{w_k\}$ and correction $\delta^* = \sum_k w_k \epsilon_k$ \hfill\Comment{Eq.~\ref{eq:mppi}}
\Else
  \State $\delta^* \leftarrow 0$ \hfill\Comment{No conflict; use consensus directly}
\EndIf
\State $\delta_{\text{opt}} \leftarrow \bar{\delta} + \delta^*$
\State \Return $p_{\text{spare}} = \sigma(\delta_{\text{opt}} / T_{\text{dec}})$ \hfill\Comment{Eq.~\ref{eq:final}}
\end{algorithmic}
\end{algorithm}

\begin{table}[t]
\caption{SWA-MPPI hyperparameter configuration. Prospect Theory parameters follow
\citet{kahneman1979prospect}; others were set via preliminary validation on a held-out
synthetic dataset and not tuned on MultiTP.}
\label{tab:hyperparams}
\centering
\begin{tabular}{llll}
\toprule
Parameter                             & Symbol                     & Value   & Source \\
\midrule
Number of persona agents              & $N$                        & 4       & Design \\
MPPI sample size                      & $K$                        & 128     & Design \\
Perturbation std.\ dev.\              & $\sigma$                   & 0.3     & Validation \\
Cooperation weight                    & $\lambda_{\text{coop}}$    & 0.7     & Validation \\
KL penalty weight                     & $\alpha_{\text{KL}}$       & 0.05    & Validation \\
PT gain curvature                     & $\alpha$                   & 0.88    & \citet{kahneman1979prospect} \\
PT loss curvature                     & $\beta$                    & 0.88    & \citet{kahneman1979prospect} \\
PT loss aversion                      & $\kappa$                   & 2.25    & \citet{kahneman1979prospect} \\
Decision temperature                  & $T_{\text{dec}}$           & 0.5     & Validation \\
Target trigger rate                   & $\rho_{\text{target}}$     & 0.35    & Design \\
Calibration set size                  & $n_{\text{cal}}$           & 50      & Design \\
Logit temp.\ (Species)                & $T_{\text{cat}}$           & 4.0     & Category \\
Logit temp.\ (Gender)                 & $T_{\text{cat}}$           & 3.5     & Category \\
Logit temp.\ (Age / Fitness / SV / U) & $T_{\text{cat}}$           & 1.5     & Category \\
\bottomrule
\end{tabular}
\end{table}

% ============================================================
\section{Experimental Setup}
\label{sec:experiments}
% ============================================================

\subsection{Dataset and Evaluation Protocol}
\label{sec:exp:data}

We evaluate on the MultiTP benchmark~\citep{jin2025multitp}, which provides trolley-problem
vignettes in 107~languages grounded in the Moral Machine dataset~\citep{awad2018moral}.
Each vignette asks an LLM to choose which of two groups a self-driving vehicle with brake
failure should spare. Groups vary across six moral dimensions (Table~\ref{tab:categories}).

For each of 15 target countries, we load vignettes in the country's \emph{native language}
and apply the following pipeline:
\begin{enumerate}
  \item \textbf{Deduplication}: retain only \texttt{paraphrase=0} vignettes (original phrasing).
  \item \textbf{Quality filtering}: remove Utilitarianism scenarios with equal group sizes on
        both sides (trivially balanced, non-discriminative).
  \item \textbf{Category capping}: limit each dimension to at most 80 scenarios to prevent
        high-frequency categories from dominating alignment metrics.
  \item \textbf{Up-sampling}: dimensions with fewer than 36 scenarios (Species and
        Utilitarianism in several languages) are augmented by \emph{random oversampling with
        replacement} of existing scenarios in that category. No paraphrasing or synthetic
        generation is performed; augmented items are exact duplicates of original vignettes
        drawn uniformly at random. This yields 36 scenarios minimum per dimension to ensure
        stable AMCE estimates.
  \item \textbf{Side balancing}: final pool is balanced to 50\% preferred-on-right by
        randomly flipping scenario group orderings as needed.
\end{enumerate}
This pipeline yields between 310 and 500 scenarios per country (342 for English/USA).
Side balance is enforced at 50\% preferred-on-right.

\textbf{Calibration/evaluation split.}
The adaptive $\tau^{(c)}$ calibration (Section~\ref{sec:method:tau}) uses $n_{\text{cal}}=50$
scenarios sampled from the scenario pool \emph{prior} to the main evaluation run.
Critically, the calibration accesses only the \emph{inter-agent variance} statistic
$\nu = \operatorname{Var}(\delta_1,\ldots,\delta_N)$ for each scenario---it never observes
the human ground-truth AMCEs or the final model alignment scores. As a result,
$\tau$ calibration is purely unsupervised with respect to the alignment objective and
cannot inflate AMCE metrics through information leakage. The same 50 scenarios are also
included in the main evaluation pool; this is valid because their variance statistics
(used only for $\tau$ estimation) carry no information about human preference values.

Ground truth AMCEs come from \texttt{country\_specific\_ACME.csv}~\citep{jin2025multitp},
aggregating 40M+ human Moral Machine responses.

\begin{table}[t]
\caption{The six moral dimensions in MultiTP with their default preferred choice
and the per-category logit temperature used in SWA-MPPI.}
\label{tab:categories}
\centering
\begin{tabular}{lllc}
\toprule
Dimension      & Preferred choice         & Contra choice       & Logit $T_{\text{cat}}$ \\
\midrule
Species        & Human                    & Animal (dog/cat)    & 4.0 \\
Gender         & Female                   & Male                & 3.5 \\
Age            & Young                    & Elderly             & 1.5 \\
Fitness        & Fit (athlete)            & Unfit (overweight)  & 1.5 \\
Social Value   & High (executive)         & Low (homeless)      & 1.5 \\
Utilitarianism & More lives               & Fewer lives         & 1.5 \\
\bottomrule
\end{tabular}
\end{table}

\subsection{Models}
\label{sec:exp:models}

We evaluate six LLMs spanning a range of sizes and families:

\begin{itemize}
  \item \textbf{Qwen2.5-72B-Instruct}~\citep{qwen2025technical}: frontier 72B instruction model (GPTQ Int8).
  \item \textbf{Meta-Llama-3.1-70B-Instruct}~\citep{grattafiori2024llama}: frontier 70B instruction model (FP8).
  \item \textbf{Qwen2.5-7B-Instruct} (4-bit quantisation): smaller 7B variant.
  \item \textbf{Meta-Llama-3.1-8B-Instruct} (4-bit quantisation): smaller 8B variant.
  \item \textbf{Mistral-Large-Instruct-2407} (126B, 4-bit quantisation)~\citep{mistral2024}: large mixture model.
  \item \textbf{Qwen2.5-32B-Instruct} (4-bit quantisation): mid-size variant.
\end{itemize}

All models are loaded on a single NVIDIA H100 80\,GB GPU via Unsloth~\citep{unsloth2024}.
All generation uses greedy decoding (temperature 0) and a fixed random seed (42).

\textbf{Cross-lingual decision token elicitation.}
Although scenarios are presented in each country's native language, the decision format
is always English. Every prompt ends with the standardised English suffix:
\begin{quote}
\small\ttfamily
Which group should be spared? Answer with LEFT or RIGHT only.
\end{quote}
The answer tokens \texttt{LEFT} and \texttt{RIGHT} are single English words that
appear as unambiguous, single-token entries in the vocabulary of all evaluated models.
We verify the token ID mapping for each model at initialisation time by encoding the
uppercase strings \texttt{"LEFT"} and \texttt{"RIGHT"} with \texttt{add\_special\_tokens=False}
and asserting that the encoding maps to exactly one token each. For Qwen2.5 models the
verified IDs are \texttt{LEFT}=22397, \texttt{RIGHT}=11572; for Llama-3.1 models
\texttt{LEFT}=31266, \texttt{RIGHT}=9268. Logits for these two token positions are then
extracted from the last-position output tensor; no generation sampling is performed.
This forced-choice setup bypasses potential refusals and off-format answers: the model
never completes the sequence---only its next-token probability mass on \texttt{LEFT}
vs.\ \texttt{RIGHT} is used. The English instruction suffix is appended to the
native-language scenario body with no separator, following the same protocol as
\citet{jin2025multitp}.

\subsection{Baselines}
\label{sec:exp:baselines}

We compare SWA-MPPI against five baselines evaluated on Llama-3.1-70B:

\begin{enumerate}
  \item \textbf{Vanilla LLM}: Single forward pass with a generic system prompt
        (\emph{``You are a helpful assistant''}); no personas or MPPI.
  \item \textbf{Country-Tailored Prompt (B1)}: System prompt that names the target country
        and asks the model to reason from that country's cultural perspective.
  \item \textbf{WVS Profile Prompting (B2)}: System prompt embedding key WVS Wave~7
        statistics for the target country (population-level averages, no age segmentation).
  \item \textbf{PRISM-Style Prompting (B3)}~\citep{kirk2024prism}: System prompt following
        the PRISM benchmark's instructions for eliciting diverse values.
  \item \textbf{Activation Steering}~\citep{arditi2025refusal}: A representative inference-time
        activation-level intervention using steering vectors computed from culturally
        contrastive prompt pairs.
\end{enumerate}

\subsection{Evaluation Metrics}
\label{sec:exp:metrics}

For each country, we compute model marginal preference rates (MPR; Eq.~\ref{eq:mpr})
as proxies for AMCEs following \citet{jin2025multitp}, and compare them to
human AMCEs using:
\begin{itemize}
  \item \textbf{JSD} (primary): Jensen-Shannon Distance $\in [0,1]$ between normalised
        model and human AMCE vectors (lower is better, $\downarrow$).
  \item \textbf{Pearson $r$}: Pearson correlation between the two 6-dimensional AMCE
        vectors (higher is better, $\uparrow$).
  \item \textbf{MAE}: Mean absolute error in AMCE percentage points ($\downarrow$).
  \item \textbf{Improvement \%}: Relative JSD reduction vs.\ the vanilla baseline
        ($\text{Improv.}\% = (1 - \text{JSD}_{\text{SWA}}/\text{JSD}_{\text{Van.}})\times 100$).
\end{itemize}
All metrics are computed per-country and aggregated as unweighted means across 15 countries.

% ============================================================
\section{Results}
\label{sec:results}
% ============================================================

\subsection{Cross-Model Comparison}
\label{sec:results:models}

Table~\ref{tab:model_summary} summarises SWA-MPPI performance across all six models.
The two frontier instruction-tuned 70B-class models benefit most: Qwen2.5-72B achieves
a mean JSD improvement of \textbf{37.6\%} and Llama-3.1-70B achieves \textbf{34.7\%},
with Pearson~$r$ rising from $0.548 \to 0.696$ and $-0.026 \to 0.384$ respectively.
Smaller 7B and 8B models show more modest but consistent improvements (+18--23\%).
In contrast, Mistral-Large-126B (+5.8\%) and Qwen2.5-32B ($-0.2\%$) benefit minimally.

We attribute the saturation of larger models to their higher baseline alignment: these
models already exhibit strong moral reasoning tendencies (Mistral vanilla Pearson $r=0.658$,
Qwen-32B $r=0.625$), leaving less room for persona-driven correction. Moreover, the
logit distributions of highly instruction-tuned models may already approximate cultural
norms better, reducing the signal available to the MPPI step.

\begin{table}[t]
\caption{Cross-model comparison of SWA-MPPI. Van.\ = vanilla baseline; SWA = SWA-MPPI.
Pearson $r$ shown for the SWA-MPPI model only. All metrics are means across 15 countries.}
\label{tab:model_summary}
\centering
\setlength{\tabcolsep}{6pt}
\begin{tabular}{lcccccc}
\toprule
\multirow{2}{*}{Model} & \multirow{2}{*}{Size} &
  \multicolumn{2}{c}{Mean JSD $\downarrow$} &
  \multirow{2}{*}{\makecell{Improv.\\ (\%)}} &
  \multicolumn{2}{c}{Mean Pearson $r$ $\uparrow$} \\
\cmidrule(lr){3-4}\cmidrule(lr){6-7}
 & & Van. & SWA & & Van. & SWA \\
\midrule
Qwen2.5-72B-Instruct      & 72B  & 0.0859 & 0.0536 & \gain{\textbf{+37.6\%}} & 0.548 & \textbf{0.696} \\
Llama-3.1-70B-Instruct    & 70B  & 0.0937 & 0.0612 & \gain{+34.7\%}          & $-$0.026 & 0.384 \\
Llama-3.1-8B-Instruct     & 8B   & 0.0533 & 0.0436 & \gain{+22.6\%}          & $-$0.247 & 0.271 \\
Qwen2.5-7B-Instruct       & 7B   & 0.0522 & 0.0426 & \gain{+18.4\%}          & 0.252  & 0.402 \\
Mistral-Large-2407        & 126B & 0.0457 & 0.0430 & \gain{+5.8\%}           & 0.658  & 0.550 \\
Qwen2.5-32B-Instruct      & 32B  & 0.0632 & 0.0633 & \loss{$-$0.2\%}         & 0.625  & 0.550 \\
\bottomrule
\end{tabular}
\end{table}

\subsection{Baseline Comparison on Llama-3.1-70B}
\label{sec:results:baselines}

Table~\ref{tab:baselines} compares SWA-MPPI against all baselines on Llama-3.1-70B across
15 countries. SWA-MPPI achieves the lowest mean JSD (0.0612) and highest mean Pearson
$r$ (0.384), outperforming all baselines. The country-tailored prompt (B1) is the
strongest alternative (mean JSD 0.0680, $+11.1\%$ gap versus SWA-MPPI). PRISM-style
prompting (B3) and WVS profile prompting (B2) achieve JSD values of 0.0699 and 0.0737
respectively. Activation steering performs \emph{below} the vanilla baseline
(JSD 0.0877), suggesting that steering vectors computed from contrastive prompts
interfere negatively with moral reasoning on this benchmark. SWA-MPPI's advantage over
B1 ($\Delta$JSD = $-10.0\%$) is particularly notable as B1 also uses country-level
cultural context---the gap demonstrates the additional value of the multi-agent MPPI
optimisation beyond simple country-level prompting.

\begin{table}[t]
\caption{Baseline comparison on Meta-Llama-3.1-70B-Instruct across 15 countries.
$\Delta$JSD (vs.\ SWA-MPPI) is the gap in mean JSD relative to our method; positive
values indicate the method is worse than SWA-MPPI.}
\label{tab:baselines}
\centering
\begin{tabular}{lcccc}
\toprule
Method & Mean JSD $\downarrow$ & Mean Pearson $r$ $\uparrow$ & $\Delta$JSD vs SWA & Gap \\
\midrule
\textbf{SWA-MPPI v3 (ours)} & \textbf{0.0612} & \textbf{0.384} & 0.000 & --- \\
Country-Tailored Prompt (B1) & 0.0680 & 0.233 & +0.0068 & +11.1\% \\
PRISM-Style Prompting (B3)   & 0.0699 & 0.148 & +0.0087 & +12.5\% \\
WVS Profile Prompting (B2)   & 0.0737 & 0.155 & +0.0125 & +16.9\% \\
Activation Steering          & 0.0877 & $-$0.208 & +0.0265 & +30.2\% \\
Vanilla Llama-3.1-70B        & 0.0937 & $-$0.026 & +0.0325 & +34.7\% \\
\bottomrule
\end{tabular}
\end{table}

\subsection{Per-Country Results: Qwen2.5-72B}
\label{sec:results:percountry}

Table~\ref{tab:results_qwen} presents detailed per-country results for Qwen2.5-72B, the
best-performing model. SWA-MPPI improves JSD in 13 of 15 countries, with the largest
gains in Australia ($+75.4\%$), Great Britain ($+72.2\%$), and the USA ($+71.1\%$).
These Anglosphere countries present an apparent paradox: despite LLMs being trained
predominantly on English data, their vanilla moral alignment with Anglosphere human
preferences is poor---vanilla JSD values of $\approx 0.10$--$0.11$ are among the highest
in our test set. This is consistent with \citet{jin2025multitp}'s finding that RLHF
alignment introduces specific moral biases (e.g., extreme protection of animals, strong
gender-neutral stances) that diverge from average human preferences in these societies.
SWA-MPPI's WVS-derived utilitarian and cohort personas successfully counterbalance
these biases.

The two countries where SWA-MPPI does not improve---Saudi Arabia ($-2.3\%$ JSD) and
Brazil ($+6.7\%$ JSD regression)---share a characteristic: limited WVS Wave~7 coverage
relative to the complexity of local moral norms. Saudi Arabia's WVS data cover a
restricted demographic sample, and Brazil shows high intra-national moral heterogeneity
that the 4-persona ensemble cannot fully capture.

\begin{table}[t]
\caption{Per-country alignment results for \textbf{Qwen2.5-72B-Instruct}.
Van.\ = Vanilla; SWA = SWA-MPPI v3. JSD\% = relative JSD reduction
(positive = improved). Best per-country SWA JSD values in \textbf{bold}.}
\label{tab:results_qwen}
\centering
\setlength{\tabcolsep}{4.5pt}
\footnotesize
\begin{tabular}{llcccccc}
\toprule
\multirow{2}{*}{Country} & \multirow{2}{*}{Lang.} &
  \multicolumn{2}{c}{JSD $\downarrow$} &
  \multicolumn{2}{c}{Pearson $r$ $\uparrow$} &
  \multicolumn{2}{c}{JSD Improv.} \\
\cmidrule(lr){3-4}\cmidrule(lr){5-6}
 & & Van. & SWA & Van. & SWA & ($\Delta$JSD) & (\%) \\
\midrule
USA         & en & 0.1066 & \textbf{0.0308} & 0.400 & 0.822 & $-$0.076 & \gain{+71.1\%} \\
Germany     & de & 0.0603 & 0.0603          & 0.845 & 0.747 & ~+0.000  & ~~+0.0\% \\
China       & zh & 0.0513 & 0.0371          & 0.851 & 0.762 & $-$0.014 & \gain{+27.8\%} \\
Japan       & ja & 0.0986 & 0.0553          & 0.634 & 0.786 & $-$0.043 & \gain{+43.9\%} \\
Brazil      & pt & 0.0754 & 0.0703          & 0.349 & 0.529 & $-$0.005 & \gain{+6.7\%}  \\
Saudi Arabia& ar & 0.0778 & 0.0796          & 0.443 & 0.527 & +0.002   & \loss{$-$2.3\%}\\
Vietnam     & vi & 0.1302 & 0.0783          & 0.569 & 0.456 & $-$0.052 & \gain{+39.9\%} \\
France      & fr & 0.0789 & 0.0479          & 0.617 & 0.838 & $-$0.031 & \gain{+39.4\%} \\
India       & hi & 0.0600 & 0.0561          & 0.670 & 0.739 & $-$0.004 & \gain{+6.5\%}  \\
South Korea & ko & 0.0720 & 0.0513          & 0.355 & 0.645 & $-$0.021 & \gain{+28.8\%} \\
Great Britain& en& 0.1057 & \textbf{0.0294} & 0.419 & 0.800 & $-$0.076 & \gain{+72.2\%} \\
Russia      & ru & 0.0800 & 0.0629          & 0.623 & 0.784 & $-$0.017 & \gain{+21.4\%} \\
Mexico      & es & 0.0760 & 0.0534          & 0.605 & 0.545 & $-$0.023 & \gain{+29.8\%} \\
Nigeria     & en & 0.1103 & 0.0649          & 0.407 & 0.615 & $-$0.046 & \gain{+41.2\%} \\
Australia   & en & 0.1047 & \textbf{0.0258} & 0.438 & 0.841 & $-$0.079 & \gain{+75.4\%} \\
\midrule
\textbf{Mean} & & \textbf{0.0859} & \textbf{0.0536} & 0.548 & 0.696 & $-$0.032 & \gain{\textbf{+37.6\%}} \\
\bottomrule
\end{tabular}
\end{table}

% ============================================================
\section{Ablation Study}
\label{sec:ablation}
% ============================================================

We conduct a comprehensive ablation study on Qwen2.5-72B evaluated on the USA with
the full 342-scenario dataset (after deduplication and balancing), isolating the
contribution of six components: the MPPI step, the persona ensemble, the utilitarian
persona specifically, the MPPI trigger strategy, the utility function, and positional
debiasing. Table~\ref{tab:ablation} reports all results; each configuration is compared
to the full SWA-MPPI system (JSD = 0.0425, Pearson $r$ = 0.724).

\paragraph{(1) No-MPPI (SWA consensus only).}
Removing the MPPI step and using the raw persona consensus $\bar{\delta}$ directly
raises JSD from 0.0425 to 0.0493 ($+15.9\%$) and reduces Pearson~$r$ by 0.074 points
(0.724 $\to$ 0.650). This confirms that the MPPI optimisation step adds meaningful
refinement beyond the consensus signal, particularly by resolving high-variance decisions
in a more culturally representative direction.

\paragraph{(2) Always-on MPPI vs.\ adaptive trigger.}
Running MPPI on \emph{every} scenario (always-on) achieves JSD = 0.0422 vs.\
JSD = 0.0425 for the adaptive trigger---a negligible difference of 0.0003.
Notably, for USA with Qwen2.5-72B, the adaptive trigger fires at 100\% (all scenarios
exceed $\tau^{(c)}$), making both strategies equivalent in practice.
This indicates that all USA scenarios present genuine inter-persona disagreement,
and the adaptive mechanism provides computational savings primarily in countries
where within-culture moral consensus is stronger. The adaptive strategy is preferred
as it gracefully degrades to direct consensus on low-disagreement scenarios without
sacrificing alignment quality.

\paragraph{(3) Utility function: Prospect Theory vs.\ alternatives.}
We compare three utility functions in the MPPI objective.
Replacing Prospect Theory with a \emph{quadratic} utility $v(x) = x^2 \cdot \text{sign}(x)$
increases JSD slightly to 0.0429 ($+0.9\%$) while retaining loss-sign asymmetry.
Replacing it with a \emph{linear} utility $v(x) = x$ raises JSD to 0.0444 ($+4.5\%$)
and reduces Pearson~$r$ by 0.018 points. The ordering PT $>$ quadratic $>$ linear
confirms that both the asymmetric curvature and the loss-aversion term of Prospect
Theory contribute to alignment quality, consistent with the hypothesis that human
moral preferences exhibit asymmetric gain-loss sensitivity~\citep{kleiman2017learning}.

\paragraph{(4) Removing the utilitarian persona.}
Dropping only the utilitarian agent (retaining three WVS cohort agents, $N = 3$)
raises JSD dramatically from 0.0425 to 0.0610 ($+43.5\%$) and collapses Pearson~$r$
from 0.725 to 0.441 ($-38.4$ points), with Spearman~$\rho$ falling from 1.000
to 0.371. The Utilitarianism dimension shows the largest per-category impact:
model preference for ``more lives'' drops from 61.2\% to 48.0\% (human: 76.6\%),
a 28.6 percentage-point gap. This reveals that the utilitarian agent is the
\emph{single most critical persona}: it anchors the ensemble's AMCE profile along
the numerosity dimension and provides the strongest corrective signal across categories.

\paragraph{(5) Positional debiasing.}
Removing the two-pass debiasing step (using a single forward pass only) raises JSD
from 0.0425 to 0.0476 ($+12.0\%$). The mean positional bias---the average absolute
difference in decision gap between original and swapped orderings---is 0.193 logit
units (median 0.093), confirming substantial recency effects. The most affected
dimension is Species: model preference for ``Human'' rises from 85.8\% (debiased)
to 93.5\% (biased), a 7.7 percentage-point inflation caused by the systematic tendency
to favour entities placed later in the prompt. The debiased configuration brings
Species closer to the human AMCE of 79.2\%.

\begin{table}[t]
\caption{Comprehensive ablation on Qwen2.5-72B-Instruct, USA (342 scenarios).
All rows compared to \textbf{Full SWA-MPPI} (adaptive trigger, PT utility, two-pass debiasing).
$\Delta$ columns are signed differences relative to the full system.}
\label{tab:ablation}
\centering
\setlength{\tabcolsep}{4pt}
\small
\begin{tabular}{llccccc}
\toprule
\# & Configuration & JSD $\downarrow$ & $\Delta$JSD & $r$ $\uparrow$ & $\Delta r$ & MAE $\downarrow$ \\
\midrule
  & \textbf{Full SWA-MPPI (adaptive, PT, 2-pass)} & \textbf{0.0425} & --- & \textbf{0.724} & --- & \textbf{9.77} \\
\midrule
1 & No-MPPI (consensus only)          & 0.0493 & \loss{+0.0068} & 0.650 & \loss{$-$0.074} & 11.02 \\
2a & Always-on MPPI                   & 0.0422 & \gain{$-$0.0003} & 0.729 & \gain{+0.005} & 9.74 \\
2b & Adaptive trigger (full system)   & 0.0425 & ---            & 0.728 & +0.004         & 9.81 \\
3a & Utility: quadratic               & 0.0429 & \loss{+0.0004} & 0.725 & +0.001         & 9.68 \\
3b & Utility: linear (no PT)          & 0.0444 & \loss{+0.0019} & 0.706 & \loss{$-$0.018} & 10.12 \\
4  & Without utilitarian persona ($N=3$) & 0.0610 & \loss{+0.0185} & 0.441 & \loss{$-$0.283} & 11.09 \\
5  & Single-pass (no debiasing)       & 0.0476 & \loss{+0.0051} & 0.743 & +0.019         & 9.67 \\
\bottomrule
\end{tabular}
\end{table}

\noindent\textbf{Summary.}
The ablation results establish a clear hierarchy of component importance:
(i)~the \emph{utilitarian persona} is most critical (JSD $+43.5\%$ without it),
followed by (ii)~the \emph{MPPI step} (JSD $+15.9\%$ without it),
(iii)~\emph{positional debiasing} (JSD $+12.0\%$ without it),
and (iv)~the \emph{Prospect Theory utility function} (JSD $+4.5\%$ with linear replacement).
The always-on vs.\ adaptive trigger distinction is negligible in the high-disagreement
USA setting, with adaptive offering practical latency benefits in lower-disagreement
countries.

% ============================================================
\section{Analysis and Discussion}
\label{sec:analysis}
% ============================================================

\subsection{Adaptive MPPI Trigger Mechanism}
\label{sec:analysis:trigger}

The adaptive tau calibration (Eq.~\ref{eq:tau}) is designed to target $\rho_{\text{target}}=35\%$
activation. For the USA with Qwen2.5-72B, however, the calibrated $\tau^{(c)}$ results
in a 100\% trigger rate across all 342 evaluation scenarios.
This is not a calibration failure but an empirical finding: every USA scenario presents
genuine inter-persona disagreement above the threshold, indicating that the six moral
dimensions all involve culturally contested trade-offs in the American context.
Practically, this makes the adaptive and always-on strategies equivalent for this setting
(Table~\ref{tab:ablation}, rows 2a--2b), and the ablation result (JSD difference of only 0.0003)
confirms no performance loss from always triggering.

The adaptive mechanism is nonetheless valuable for other countries: cultures with stronger
internal moral consensus (e.g., countries with more homogeneous WVS profiles) will exhibit
lower trigger rates and thus skip the MPPI step entirely for straightforward scenarios.
This graceful degradation ensures SWA-MPPI adds computational overhead only where
cultural complexity warrants it.

Importantly, the \emph{flip rate}---the fraction of triggered scenarios where MPPI
reverses the vanilla decision direction---is consistently low (mean 6.0\% across all
evaluated countries). This indicates that MPPI primarily \emph{modulates decision confidence}
rather than reversing choices wholesale. Strong vanilla commitments are preserved;
ambiguous decisions are resolved in a culturally grounded direction. This is a desirable
property: the framework avoids overcorrection and catastrophic interference with
well-calibrated model outputs.

\subsection{Latency Analysis}
\label{sec:analysis:latency}

The mean per-scenario latency for SWA-MPPI is approximately 730\,ms on Qwen2.5-72B
(H100 80\,GB), compared to 160\,ms for the vanilla single-pass baseline---a factor of
$\approx 4.6\times$ overhead. This overhead stems from two sources: (i) the batched
multi-agent forward pass ($N+1=5$ sequences per pass), and (ii) the positional debiasing
second pass. The MPPI optimisation loop ($K=128$ perturbations) operates entirely on
cached scalar logit values and contributes negligibly to wall-clock time. At 500 scenarios
per country and 15 countries, total evaluation time per model is approximately 91 minutes.

\subsection{Failure Cases and Limitations}
\label{sec:analysis:limitations}

\paragraph{Models with high baseline alignment.}
Mistral-Large-126B and Qwen2.5-32B exhibit vanilla Pearson~$r$ values of $0.658$ and
$0.625$ respectively, indicating that these models already encode substantial cultural
sensitivity---possibly due to multilingual training corpora and more diverse RLHF
annotations. For these models, the persona ensemble introduces noise without sufficient
corrective signal, and SWA-MPPI provides marginal or slightly negative improvement.
Future work should investigate dynamic selection of whether to apply SWA-MPPI based
on a preliminary estimate of baseline alignment quality.

\paragraph{Limited WVS coverage.}
13 of 15 target countries have WVS Wave~7 coverage; Saudi Arabia and Vietnam rely on
partial data. Countries outside WVS coverage (many sub-Saharan African and Pacific
Island nations) would require manually authored personas, reducing empirical grounding
and risking stereotyping~\citep{deshpande2023toxicity}.

\paragraph{Binary decision format.}
The MultiTP benchmark presents strictly binary forced choices. SWA-MPPI is designed for
this format; extension to multi-option or continuous preference elicitation is left for
future work. In practice, many moral decisions are not binary, and the logit-difference
formulation ($\delta = z_r - z_\ell$) would need to be generalised to a softmax over
multiple options.

\paragraph{Language--culture conflation.}
Following \citet{jin2025multitp}, we use one language per country, potentially conflating
linguistic and cultural effects. Countries with multiple dominant languages (India,
Nigeria, Switzerland) are each represented by a single language variety, missing
intra-national diversity. Disentangling language from culture remains an important
open problem.

\paragraph{Model quantisation.}
All models are loaded in 4-bit or Int8 quantisation. While quantisation loss is typically
small for 70B+ models, it may introduce systematic logit artefacts that affect the
numerical stability of the MPPI optimisation step.

\subsection{Broader Impact}
\label{sec:analysis:impact}

SWA-MPPI provides a principled mechanism for adapting LLM moral behaviour to diverse
cultural contexts without modifying model weights. This has both positive and potentially
concerning implications. On the positive side, the framework could support more culturally
equitable AI assistants that serve non-WEIRD populations more faithfully. On the
concerning side, the same framework could be used to align LLMs with moral preferences
that are harmful by other ethical standards (e.g., preferences for discrimination against
specific groups). The WVS-grounding provides a degree of empirical accountability---persona
descriptions reflect surveyed population averages rather than individual author
preferences---but it does not eliminate the risk of culturally legitimised harm. We
advocate for careful human oversight when deploying culturally adaptive alignment systems,
and note that SWA-MPPI does not override safety-critical model guardrails.

% ============================================================
\section{Conclusion}
\label{sec:conclusion}
% ============================================================

We have presented SWA-MPPI, a training-free, inference-time framework for cross-cultural
moral alignment of large language models. The method constructs culturally grounded
persona agents from World Values Survey data, evaluates their logit-space consensus in
a single batched forward pass, and applies a Prospect-Theory-weighted Model Predictive
Path Integral correction whenever inter-agent disagreement signals potential cultural
misalignment. A positional debiasing step further reduces systematic left/right ordering
artefacts.

Evaluated on the MultiTP benchmark across 15 countries, 12 native languages, and 6 moral
dimensions, SWA-MPPI reduces Jensen-Shannon Distance to human moral preferences by
\textbf{37.6\%} on Qwen2.5-72B and \textbf{34.7\%} on Llama-3.1-70B, while improving
Pearson correlation with human AMCE vectors by up to $+0.41$ points. Improvements are
largest for Anglosphere countries (USA, GBR, AUS: $+71$--$75\%$) where RLHF-induced
moral biases are strongest, and for culturally distant country--language pairs
(Vietnam, Mexico) where vanilla alignment is particularly poor. Ablation studies
demonstrate that WVS-grounded personas are the critical component, MPPI adds
meaningful gain beyond the persona consensus, and the Prospect Theory utility function
provides a beneficial inductive bias.

SWA-MPPI operates entirely on frozen model weights, requires no labelled preference data,
and introduces a mean per-scenario overhead of $\approx 730$\,ms on a single H100
GPU---a practical cost for culturally aware inference at deployment scale. We hope this
work contributes a concrete step toward the pluralistic alignment vision
of~\citet{sorensen2024roadmap}: AI systems that genuinely model the diversity of human
moral preferences across cultures, rather than imposing a single, culturally unanchored
ethical stance on a diverse world.

\paragraph{Future work.}
Promising directions include: extending the persona library to cover additional WVS
variables and non-WVS countries; investigating interaction of SWA-MPPI with fine-tuned
models and multilingual base models; integrating structured chain-of-thought reasoning
into persona responses; applying the MPPI paradigm to multi-option preference scenarios;
and developing a meta-learning criterion for deciding whether to activate SWA-MPPI based
on estimated baseline alignment quality.

% ============================================================
\begin{ack}
We thank the creators of the MultiTP dataset~\citep{jin2025multitp} and the Moral Machine
experiment~\citep{awad2018moral} for making their data and benchmarks publicly available.
We thank the World Values Survey Association~\citep{haerpfer2022wvs} for providing freely
accessible cross-national survey data. Experiments were conducted on NVIDIA H100 hardware.
\end{ack}

% ============================================================
\bibliographystyle{abbrvnat}
\bibliography{references}

\end{document}
